\bibitem{Beilinson}
A. A. Beilinson,
\emph{Coherent sheaves on \(\PP^n\) and problems of linear algebra},
Functional Analysis and Its Applications \textbf{12} (1978), 214--216.
\url{https://doi.org/10.1007/BF01681436}.
原著Section 4の対角Koszul分解を確認。
\url{https://www.mathnet.ru/eng/faa2009}.
\bibitem{DDL}
O. De Deken and W. Lowen,
\emph{Abelian and derived deformations in the presence of \(\ZZ\)-generating geometric helices},
Journal of Noncommutative Geometry \textbf{5} (2011), 477--505.
Theorem 4.5; Propositions 5.5, 5.8, 5.9; Lemma 5.14.
\url{https://doi.org/10.4171/JNCG/83}.
\bibitem{Keller}
B. Keller, \emph{Deriving DG categories},
Annales scientifiques de l'École Normale Supérieure \textbf{27} (1994), 63--102.
\url{https://doi.org/10.24033/asens.1689}.
\bibitem{Mathew}
A. Mathew, \emph{The Galois group of a stable homotopy theory},
Advances in Mathematics \textbf{291} (2016), 403--541.
Proposition 3.28, Remark 3.29, Proposition 3.35.
\url{https://arxiv.org/abs/1404.2156}.
\bibitem{Orlov}
D. Orlov,
\emph{Derived categories of coherent sheaves and triangulated categories of singularities},
arXiv:math/0503632, v3 (2009), Theorem 2.5.
\url{https://arxiv.org/abs/math/0503632}.
\bibitem{CCKM}
T. Cassidy, A. Conner, E. Kirkman and W. F. Moore,
\emph{Periodic free resolutions from twisted matrix factorizations},
Journal of Algebra \textbf{455} (2016), 137--163.
Proposition 5.7, Theorem 5.8.
\url{https://doi.org/10.1016/j.jalgebra.2016.01.037}.
\bibitem{CT}
A. Căldăraru and J. Tu, \emph{Curved \(A_\infty\) algebras and Landau--Ginzburg models},
New York Journal of Mathematics \textbf{19} (2013), 305--342.
Section 5、特にdg Orlov windowのTheorem 5.8.
\url{https://nyjm.albany.edu/j/2013/19-16v.pdf}.
\bibitem{BFK}
M. Ballard, D. Favero and L. Katzarkov,
\emph{A category of kernels for equivariant factorizations and its implications for Hodge theory},
Publications Mathématiques de l'IHÉS \textbf{120} (2014), 1--111.
Theorem 5.39のKoszul cohomologyによる公式、Theorem 5.15とCorollary 5.18のequivariant外部積.
\url{https://doi.org/10.1007/s10240-013-0059-9}.
確認した公刊本文 \url{https://www.numdam.org/article/PMIHES_2014__120__1_0.pdf}.
\bibitem{LZ}
X. Lin and S. Zhang,
\emph{Serre algebra, matrix factorization and categorical Torelli theorem for hypersurfaces},
arXiv:2310.09927, v2 (2024), Theorem 1.7.
\url{https://arxiv.org/abs/2310.09927}.
\bibitem{HR}
D. Huybrechts and J. V. Rennemo,
\emph{Hochschild cohomology versus the Jacobian ring, and the Torelli theorem for cubic fourfolds},
Algebraic Geometry \textbf{6} (2019), 76--99.
Lemma 2.19およびSections 2--3.
\url{https://doi.org/10.14231/AG-2019-005}.
\bibitem{Huy}
D. Huybrechts, \emph{The K3 category of a cubic fourfold},
Compositio Mathematica \textbf{153} (2017), 586--620.
Proposition 3.12およびSection 3.5.
\url{https://doi.org/10.1112/S0010437X16008137}.
\bibitem{JLLZ}
A. Jacovskis, X. Lin, Z. Liu and S. Zhang,
\emph{Infinitesimal categorical Torelli theorems for Fano threefolds},
arXiv:2203.08187, Theorem 4.4(2).
\url{https://arxiv.org/abs/2203.08187}.
\bibitem{Lurie}
J. Lurie, \emph{Derived Algebraic Geometry X: Formal Moduli Problems},
2011. Theorems 5.1.9, 5.3.16およびProposition 5.3.21.
\url{https://www.math.ias.edu/~lurie/papers/DAG-X.pdf}.
\bibitem{BKP}
A. Blanc, L. Katzarkov and P. Pandit,
\emph{Generators in formal deformations of categories},
Compositio Mathematica \textbf{154} (2018).
本文で用いるのはarXiv版のDefinition 3.3およびSection 3.3.
\url{https://arxiv.org/abs/1705.00655}.
\bibitem{KL}
B. Keller and W. Lowen,
\emph{On Hochschild cohomology and Morita deformations},
International Mathematics Research Notices (2009).
\url{https://webusers.imj-prg.fr/~bernhard.keller/publ/KellerLowenHochDefo.pdf}.
\bibitem{EG}
P. Etingof and V. Ginzburg,
\emph{Noncommutative del Pezzo surfaces and Calabi--Yau algebras},
Journal of the European Mathematical Society \textbf{12} (2010), 1371--1416.
Sections 3.5--3.6, Theorem 3.6.2, Proposition 8.1.1(i) およびその証明。
\url{https://doi.org/10.4171/JEMS/235}.
\href{https://arxiv.org/abs/0709.3593v4}{arXiv:0709.3593v4}.
\bibitem{Walton}
C. Walton, \emph{Representation theory of three-dimensional Sklyanin algebras},
Nuclear Physics B \textbf{860} (2012), 167--185.
Theorem 2.14および式(2.15)--(2.16).
\url{https://doi.org/10.1016/j.nuclphysb.2012.02.015}.
\bibitem{CGWZ}
K. Chan, J. Gaddis, R. Won and J. J. Zhang,
\emph{Ozone groups of Artin--Schelter regular algebras satisfying a polynomial identity},
arXiv:2312.13014, v2 (2024). Introductionの非退化Sklyanin代数の記述。
\url{https://arxiv.org/abs/2312.13014}.
\bibitem{Elagin}
A. Elagin, \emph{On equivariant triangulated categories},
arXiv:1403.7027, Theorem 4.2.
\url{https://arxiv.org/abs/1403.7027}.
\bibitem{ELS}
A. Elagin, V. A. Lunts and O. M. Schn\"urer,
\emph{Smoothness of Derived Categories of Algebras},
Moscow Mathematical Journal \textbf{20} (2020), no. 2, 277--309.
確認した一次資料はarXiv:1810.07626v2 (22 March 2019),
Theorem 4.4 (Verdier局所化), Proposition 4.10とTheorem 4.15 (稠密開集合と生成子),
Theorem 4.18 (外部積生成), Theorem 5.1 (有限中心代数の有界導来圏のsmoothness).
\url{https://arxiv.org/abs/1810.07626};
\url{https://www.mathjournals.org/mmj/2020-020-002/2020-020-002-003.html}.

\bibitem{TV}
B. Toën and M. Vaquié,
\emph{Moduli of objects in dg-categories},
Annales scientifiques de l'École Normale Supérieure (4) \textbf{40} (2007), 387--444.
Theorem 3.6; Corollaries 2.13, 3.15, 3.19, 3.21; Lemma 3.25 and proof of Corollary 3.24.
\url{https://arxiv.org/abs/math/0503269};
確認した著者版 \url{https://www.math.univ-toulouse.fr/~btoen/ttt.pdf}.
\bibitem{Stacks}
The Stacks Project Authors, \emph{The Stacks Project}.
Tag 02HX (Lemma 37.12.1, Noetherian底上の無限小smoothness判定),
Tag 055U (Lemma 37.38.6, smooth morphismのétale局所section).
\url{https://stacks.math.columbia.edu/tag/02HX};
\url{https://stacks.math.columbia.edu/tag/055U}.
参照日2026年11月4日。
\bibitem{Pirozhkov}
D. Pirozhkov, \emph{A categorical Torelli theorem for hypersurfaces},
Bulletin of the London Mathematical Society \textbf{56} (2024), 3075--3089.
\url{https://doi.org/10.1112/blms.13117}.
公刊版のabstractで、cubicに対する仮定 \(n>3\) を確認。
本稿の全次元の一意性には \cite[Theorem 1.7]{LZ} を用いる。
\bibitem{HuyUpdate}
D. Huybrechts, \emph{The K3 category of a cubic fourfold: an update},
Beiträge zur Algebra und Geometrie / Contributions to Algebra and Geometry (2025).
Section 9: Unknown deformations.
\url{https://doi.org/10.1007/s13366-025-00788-9}.

\bibitem{ChanMulti}
D. Chan, \emph{Twisted multi-homogeneous coordinate rings},
Journal of Algebra \textbf{223} (2000), 438--456.
Theorem 3.4; proof of Theorem 5.2, Claim 5.3 and Lemmas 5.4--5.6;
Corollaries 5.8--5.9.
\url{https://doi.org/10.1006/jabr.1999.7938}.
確認した著者版 \url{https://web.maths.unsw.edu.au/~danielch/paper/tbcr.ps}, pp.8--9,14--17.
\bibitem{SmithTate}
S. P. Smith and J. Tate,
\emph{The center of the 3-dimensional and 4-dimensional Sklyanin algebras},
K-Theory \textbf{8} (1994), 19--63.
Corollary 2.3 (p.~26); Theorem 3.7(c) (p.~38).
\url{https://doi.org/10.1007/BF00962089};
\url{https://sites.math.washington.edu/~smith/Research/tate.pdf}.
\bibitem{KeelerMulti}
D. S. Keeler, \emph{Noncommutative ampleness for multiple divisors},
Journal of Algebra \textbf{265} (2003), no. 1, 299--311.
Theorem 2.7, Corollary 2.8 (任意有限個のbimoduleのNC-ampleness);
Hypothesis 3.1, Theorem 3.3, Corollary 3.5 (二因子のnoetherian性).
番号はarXiv一次TeXソースで確認し、刊行版との番号一致は未照合。
\url{https://doi.org/10.1016/S0021-8693(03)00126-1};
\url{https://arxiv.org/abs/math/0210417};
\url{https://arxiv.org/e-print/math/0210417}.
\bibitem{WWY3}
C. Walton, X. Wang and M. Yakimov,
\emph{Poisson geometry of PI three-dimensional Sklyanin algebras},
Proceedings of the London Mathematical Society \textbf{118} (2019), 1471--1500.
Theorem 1.3(1),(2a), verified in arXiv:1704.04975v3, p.~3.
\url{https://doi.org/10.1112/plms.12220};
\url{https://arxiv.org/abs/1704.04975}.
\bibitem{ToenFinitude}
B. Toën, \emph{Finitude homotopique des dg-algèbres propres et lisses},
Proceedings of the London Mathematical Society (3) \textbf{98} (2009), 217--240.
Lemmas 3.13--3.14. Lemma 3.14の結論は部分前層のquasi-compactnessであり、一般の開性ではない。
\url{https://doi.org/10.1112/plms/pdn027};
確認した著者版 \url{https://www.math-beb.fr/wp-content/uploads/2026/10/finitesatdgalg.pdf}.

\bibitem{SnowdenDVR}
A. Snowden, \emph{Lecture 8: Elliptic curves over DVRs},
Math 679 lecture notes, University of Michigan, 2013.
Section ``Reduction of torsion points'', good reductionでのprime-to-residue-characteristic torsionの還元の単射性。
\url{https://websites.umich.edu/~asnowden/teaching/2013/679/L08.html}.
参照日2026年11月6日。
\bibitem{HMY}
J.-W. He, X.-C. Ma and Y. Ye, \emph{Generalized Kn\"orrer's Periodicity Theorem},
arXiv:2107.06438v2 (16 August 2021).
対象は中心二次元の商であり、本稿の三次の場合には直接適用しない。
\url{https://arxiv.org/abs/2107.06438}.
